% Theorem numbers for Efimov_Large_Categories refer to arXiv:2405.12169v4.
	\chapter{Examples of dualizable categories}
	\label[chapter]{chap:Examples}

	% Source: Talk 3 (Max Glaser, 27.10.2026)
	In this chapter we discuss several sources of dualizable (and, more generally, compactly assembled) categories:
	\begin{itemize}
		\item Kernels of homological epimorphisms of ring spectra, and in particular almost ring theory (\Cref{sec:Homological_Epimorphisms,sec:Almost_Ring_Theory}).
		\item Continuous posets, as a class of non-stable compactly assembled categories (\Cref{sec:Continuous_Posets}).
		\item Gluing constructions and extensions, including a nondualizable extension of compactly generated categories (\Cref{sec:Extensions_And_Gluing}).
	\end{itemize}

	\section{Kernels of homological epimorphisms}
	\label[section]{sec:Homological_Epimorphisms}

	% Source: Clausen, Lecture 2 (18.07.2023); Mathew, Lecture 3 (MPIM Bonn, 10.09.2024)
	Let $A \to B$ be a map of $\bbE_1$-rings, and consider the base change functor $B \otimes_A -\colon \D(A) \to \D(B)$ between the categories of left modules, with right adjoint the restriction functor $\res\colon \D(B) \to \D(A)$. Recall from \Cref{def:Homological_Epimorphism} that $A \to B$ (or rather $\Perf(A) \to \Perf(B)$) is a \emph{homological epimorphism} if $B \otimes_A -$ is a localization, i.e.\ if $\res$ is fully faithful. Since $\res$ preserves colimits, this is equivalent to the counit $B \otimes_A \res(N) \to N$ being an equivalence for all $N \in \D(B)$, and since both sides preserve colimits in $N$ and $\D(B)$ is generated by $B$, this is in turn equivalent to the multiplication map
	\[
		B \otimes_A B \to B
	\]
	being an equivalence. In the language of Clausen, $A \to B$ is then called \emph{idempotent}.

	\begin{proposition}
		\label[proposition]{prop:Kernel_Homological_Epi}
		Let $A \to B$ be a homological epimorphism of $\bbE_1$-rings, and put
		\[
			\D(A \to B) := \ker\big(B \otimes_A -\colon \D(A) \to \D(B)\big).
		\]
		Then $\D(A \to B)$ is dualizable, and the sequence $\D(A \to B) \hookrightarrow \D(A) \to \D(B)$ is a strongly continuous localization sequence. In particular,
		\[
			K^{\cont}(\D(A \to B)) \simeq \fib\big(K(A) \to K(B)\big).
		\]
	\end{proposition}
	\prooffor{prop:Kernel_Homological_Epi}
	\begin{proof}
		The functor $B \otimes_A -$ is a localization by assumption, and it is strongly continuous since its right adjoint $\res$ preserves colimits. Hence $\D(A \to B) \hookrightarrow \D(A) \to \D(B)$ is a strongly continuous localization sequence (\Cref{lem:Strong_Continuity_Localization_Sequence}), and $\D(A \to B)$ is dualizable by \Cref{prop:Localization_Sequence_Dualizable}. The final statement follows from the definition of $K^{\cont}$, see \Cref{sec:Continuous_K_Theory_Intro}.
	\end{proof}

	\begin{proposition}
		\label[proposition]{prop:Every_Dualizable_Kernel}
		Every dualizable category is of the form $\D(E \to E')$ for some homological epimorphism $E \to E'$ of $\bbE_1$-rings.
	\end{proposition}
	\prooffor{prop:Every_Dualizable_Kernel}
	\begin{proof}
		Let $C$ be dualizable. Choose a set $S$ of representatives of the isomorphism classes of $C^{\omega_1}$, a regular cardinal $\kappa > |S|$, and put $G := \bigoplus_{x \in S} x$. As a $\kappa$-small sum of $\kappa$-compact objects, $G$ is $\kappa$-compact. Let $A \subseteq C^{\kappa}$ be the thick subcategory generated by $G$; it contains $C^{\omega_1}$, since every object of $C^{\omega_1}$ is a retract of $G$. The functor $\hat{y}\colon C \to \Ind(C^{\kappa})$ takes values in $\Ind(C^{\omega_1}) \subseteq \Ind(A)$ (\Cref{thm:Compactly_Assembled_Omega_1}, see also \cite[Corollary~2.14]{Efimov_Large_Categories}), and hence is a fully faithful left adjoint of $\colim\colon \Ind(A) \to C$. Its essential image $\hat{y}(C)$ is closed under colimits, with right orthogonal $\ker(\colim\colon \Ind(A) \to C)$. Since the latter is closed under colimits in $\Ind(A)$, the localization
		\[
			Q\colon \Ind(A) \to \Ind(A)/\hat{y}(C) \simeq \ker(\colim)
		\]
		has a colimit-preserving right adjoint and thus preserves compact objects. As $Q$ is essentially surjective and $\Ind(A)$ is compactly generated by $G$, the category $\ker(\colim)$ is compactly generated by the single object $Q(G)$. Hence $\ker(\colim) \simeq \Ind(B)$ for $B$ the thick subcategory generated by $Q(G)$, and $Q$ is $\Ind$ of a homological epimorphism $A \to B$ with kernel $\hat{y}(C) \simeq C$; this is the Calkin sequence of \Cref{thm:Lurie_Dualizable_Characterization} with $C^{\kappa}$ replaced by $A$. Finally, by Morita theory \cite[Theorem~7.1.2.1]{lurie2016HA} we have $\Ind(A) \simeq \D(E)$ for $E := \End(G)\catop$ and $\Ind(B) \simeq \D(E')$ for $E' := \End(Q(G))\catop$, under which $Q$ corresponds to $E' \otimes_E -$, where the ring map $E \to E'$ is induced by $Q$ on endomorphisms. Indeed, $Q(G)$ has the resulting $E$-action, so the two colimit-preserving functors agree on the generator together with its endomorphism action, as required by Morita theory. So $C \simeq \D(E \to E')$.
	\end{proof}
	% In Clausen's lectures this was attributed to Efimov; we have not found it stated in this form in Efimov's papers. The proof above was written with AI assistance and should be checked.

	There is a convenient criterion for a map $A \to B$ with fiber $I$ to be a homological epimorphism. Recall that a non-unital ring spectrum $I$ is called \emph{h-unital} if the multiplication $I \otimes_{\S\oplus I} I \to I$ is an isomorphism, where $\S\oplus I$ denotes the unitalization. Equivalently, the nonunital relative tensor product, defined by the semisimplicial bar construction with terms $I^{\otimes(n+2)}$, maps isomorphically to $I$; see \cite[Section~2.10 and Lemma~3.5.5]{Krause_Nikolaus_Putzstuck_Sheaves}.

	\begin{theorem}[{\cite[Section~2.10, especially Corollary~2.10.6]{Krause_Nikolaus_Putzstuck_Sheaves}}]
		\label[theorem]{thm:H_Unital}
		For a non-unital ring spectrum $I$, the following are equivalent:
		\begin{enumerate}
			\item $I$ is h-unital;
			\item the map $\S \oplus I \to \S$ is a homological epimorphism;
			\item for every fiber sequence $I \to A \to B$ in which $A \to B$ is a map of unital ring spectra, the map $A \to B$ is a homological epimorphism.
		\end{enumerate}
	\end{theorem}

	Conversely, if $A \to B$ is a homological epimorphism of connective ring spectra and $\pi_0(A) \to \pi_0(B)$ is surjective, then $\fib(A \to B)$ is h-unital \cite[Proposition~2.10.7]{Krause_Nikolaus_Putzstuck_Sheaves}. The surjectivity hypothesis cannot be omitted: the fiber of $k[t]\to k[t,t^{-1}]$ is not h-unital \cite[Example~2.10.8]{Krause_Nikolaus_Putzstuck_Sheaves}.

	\begin{example}
		\label[example]{ex:Homological_Epis}
		The following maps are homological epimorphisms, and hence give rise to dualizable categories $\D(A \to B)$:
		\begin{enumerate}
			\item A localization $A \to A[S^{-1}]$ of a commutative ring at a multiplicative subset. More generally, if $A$ is a commutative Noetherian ring, the kernel of every homological epimorphism out of $A$ is compactly generated: the associated localization of $\D(A)$ is smashing, so this follows from \cite[Theorem~3.3 and Corollary~3.4]{Neeman_Chromatic_Tower}.
			\item A surjection $A \twoheadrightarrow B$ of perfect $\F_p$-algebras: the higher $\mathrm{Tor}^A$-groups between perfect $A$-algebras vanish by \cite[Lemma~3.16]{Bhatt_Scholze_Witt_Grassmannian}, so $B \otimes^{\bL}_A B \simeq B$. Similarly, if $A$ is an integral perfectoid ring and $B$ is a perfect $\F_p$-algebra, any surjection $A \twoheadrightarrow B$ is a homological epimorphism. Indeed, the derived $p$-completed tensor product of perfectoid rings is discrete by the proof of \cite[Lemma~10.3]{Bhatt_Scholze_Prisms}. Here $B \otimes^{\bL}_A B$ is already derived $p$-complete, since it is a $B$-module, and its degree-zero homology is $B$.
			\item (Wodzicki) If $A$ is a unital C$^*$-algebra and $I \subseteq A$ a closed two-sided $*$-ideal, then $A \to A/I$ is a homological epimorphism. For $A = B(\Hh)$ and $I = \Kk(\Hh)$, this gives the category $\D(\Kk)$ from \Cref{ex:Compact_Operators_Intro}.
			\item (Almost ring theory) If $A$ is a ring and $I \subseteq A$ is a two-sided ideal with $I^2 = I$ which is flat as a left or right $A$-module, then $A \to A/I$ is a homological epimorphism. We discuss this in detail in \Cref{sec:Almost_Ring_Theory}.
		\end{enumerate}
	\end{example}
	% Source: written from the literature (Neeman, Theorem 3.3; Bhatt--Scholze, Witt vectors, Lemma 3.16, and Prisms, proof of Lemma 10.3).

	\section{Almost ring theory}
	\label[section]{sec:Almost_Ring_Theory}

	% Source: Talk 3 (Max Glaser, 27.10.2026)
	Almost ring theory (Faltings, Gabber--Ramero \cite{Gabber_Ramero_Almost}) provides one of the most natural sources of dualizable categories which are not compactly generated. Throughout this section, let $A$ be an (associative, unital) ring and let $I \subseteq A$ be a two-sided ideal such that
	\begin{itemize}
		\item $I^2 = I$, and
		\item $I$ is flat as a right $A$-module.
	\end{itemize}
	We work with left $A$-modules, and write $A\text{-}\Mod$ for the abelian category of left $A$-modules and $\D(A)$ for its derived category.

	The forgetful functor $(A/I)\text{-}\Mod \hookrightarrow A\text{-}\Mod$ is fully faithful, with essential image the $A$-modules $M$ with $IM = 0$. Since $I^2 = I$, this subcategory is closed under extensions: if $0 \to M' \to M \to M'' \to 0$ is exact with $IM' = 0 = IM''$, then $IM = I^2M \subseteq I M' = 0$. As it is clearly closed under subobjects and quotients, it is a \emph{Serre subcategory}.

	\begin{definition}
		\label[definition]{def:Almost_Modules}
		The category of \emph{$(A,I)$-almost modules} is the Serre quotient
		\[
			(A,I)^a\text{-}\Mod := A\text{-}\Mod/(A/I)\text{-}\Mod.
		\]
	\end{definition}

	\begin{example}
		\label[example]{ex:Almost_Modules}
		\leavevmode
		\begin{enumerate}
			\item If $I = A$, then $(A/I)\text{-}\Mod = 0$ and $(A,A)^a\text{-}\Mod = A\text{-}\Mod$.
			\item The motivating examples come from non-discrete valuation rings: let $A$ be a valuation ring of rank one whose value group is not discrete, e.g.\ $A = \Oo_{\C_p}$, and let $I = \mm$ be its maximal ideal. Then $\mm^2 = \mm$, since every element of $\mm$ is a product of two elements of $\mm$, and $\mm$ is flat since it is torsion-free over a valuation ring. An $A$-module is called \emph{almost zero} if it is killed by $\mm$; for instance, the residue field $A/\mm$ becomes zero in the almost category. This is the setting of Faltings' almost purity theorem in $p$-adic Hodge theory.
		\end{enumerate}
	\end{example}

	\subsection*{The derived picture}
	We have an adjoint triple
	\[
		\begin{tikzcd}[column sep=large]
			\D(A) \arrow[r, shift left=2.5, "L"] \arrow[r, shift right=2.5, "R"'] & \D(A/I) \arrow[l, hook', "\res" description]
		\end{tikzcd}
	\]
	where $\res$ denotes restriction of scalars, $L = A/I \otimes^{\bL}_A -$ and $R = \bR\Hom_A(A/I,-)$.

	\begin{lemma}
		\label[lemma]{lem:Almost_Restriction_FF}
		The functor $\res\colon \D(A/I) \to \D(A)$ is fully faithful. In particular, $L$ is a localization and $A \to A/I$ is a homological epimorphism.
	\end{lemma}
	\prooffor{lem:Almost_Restriction_FF}
	\begin{proof}
		As explained at the beginning of \Cref{sec:Homological_Epimorphisms}, it suffices to show that the multiplication map $A/I \otimes^{\bL}_A A/I \to A/I$ is an equivalence. Using the short exact sequence $0 \to I \to A \to A/I \to 0$ of right $A$-modules, we find
		\begin{align*}
			\mathrm{Tor}^A_0(A/I,A/I) &= A/I, \\
			\mathrm{Tor}^A_1(A/I,A/I) &\cong I/I^2 = 0, \\
			\mathrm{Tor}^A_k(A/I,A/I) &\cong \mathrm{Tor}^A_{k-1}(I,A/I) = 0 \quad \text{for } k \geq 2,
		\end{align*}
		where the last vanishing uses flatness of $I$ as a right $A$-module.
	\end{proof}

	\begin{lemma}
		\label[lemma]{lem:Almost_Essential_Image}
		The essential image of $\res\colon \D(A/I) \hookrightarrow \D(A)$ is the full subcategory
		\[
			\D_{(A/I)\text{-}\Mod}(A) := \{X \in \D(A) \mid IH_n(X) = 0 \text{ for all } n \in \Z\}
		\]
		of complexes whose homology lies in the Serre subcategory $(A/I)\text{-}\Mod$. Moreover, the kernel of $L$ consists of those $X \in \D(A)$ with $X \simeq I \otimes^{\bL}_A X$, and $I \otimes^{\bL}_A X \in \ker(L)$ for every $X \in \D(A)$.
	\end{lemma}
	\prooffor{lem:Almost_Essential_Image}
	\begin{proof}
		Since $\res$ is fully faithful with left adjoint $L$, its essential image consists of those $X$ for which the unit $X \to \res L(X)$ is an equivalence. Tensoring the short exact sequence $0 \to I \to A \to A/I \to 0$ of $A$-bimodules with $X$ gives a cofiber sequence
		\[
			I \otimes^{\bL}_A X \to X \to \res L(X),
		\]
		so $X$ lies in the essential image if and only if $I \otimes^{\bL}_A X \simeq 0$. Since $I$ is flat as a right module, $H_n(I \otimes^{\bL}_A X) \cong I \otimes_A H_n(X)$. Finally, for a left $A$-module $M$, we have $I \otimes_A M = 0$ if and only if $IM = 0$: the multiplication map $I \otimes_A M \to IM$ is surjective, and conversely if $IM = 0$ then $M$ is an $A/I$-module and $I \otimes_A M \cong (I \otimes_A A/I) \otimes_{A/I} M \cong I/I^2 \otimes_{A/I} M = 0$.

		For the second statement, $L(I \otimes^{\bL}_A X) \simeq (A/I \otimes^{\bL}_A I) \otimes^{\bL}_A X$, and $A/I \otimes^{\bL}_A I \simeq \fib(A/I \to A/I \otimes^{\bL}_A A/I) \simeq 0$ by \Cref{lem:Almost_Restriction_FF}. The characterization of $\ker(L)$ follows from the cofiber sequence above.
	\end{proof}

	\begin{proposition}
		\label[proposition]{prop:Almost_Derived_Category}
		There is an equivalence
		\[
			\D((A,I)^a\text{-}\Mod) \simeq \ker\big(L\colon \D(A) \to \D(A/I)\big),
		\]
		and $\ker(L) = \D(A \to A/I)$ is dualizable. More precisely, $\D((A,I)^a\text{-}\Mod) \hookrightarrow \D(A) \xrightarrow{L} \D(A/I)$ is a strongly continuous localization sequence, and $\D((A,I)^a\text{-}\Mod)$ is also the kernel of $L$ computed in $\Pr^{\dual}$.
	\end{proposition}
	\prooffor{prop:Almost_Derived_Category}
	\begin{proof}
		Since $(A/I)\text{-}\Mod \subseteq A\text{-}\Mod$ is a Serre subcategory of a Grothendieck abelian category which is closed under colimits, the derived category of the Serre quotient is the Verdier quotient
		\[
			\D((A,I)^a\text{-}\Mod) \simeq \D(A)/\D_{(A/I)\text{-}\Mod}(A).
		\]
		Indeed, the quotient functor $A\text{-}\Mod \to (A,I)^a\text{-}\Mod$ is exact and has a fully faithful right adjoint, and for such functors this identification is proved in \cite[Section~7.7]{Krause_Localization}; see also \cite[Example~2.26]{Efimov_Large_Categories}. It thus remains to identify the Verdier quotient $\D(A)/\D_{(A/I)\text{-}\Mod}(A)$ with $\ker(L)$. Consider the functor $q := I \otimes^{\bL}_A -\colon \D(A) \to \ker(L)$, which lands in $\ker(L)$ by \Cref{lem:Almost_Essential_Image}. For $K \in \ker(L)$ and $X \in \D(A)$, applying $\Map(K,-)$ to the cofiber sequence $I \otimes^{\bL}_A X \to X \to \res L(X)$ and using $\Map(K,\res L X) \simeq \Map(LK,LX) \simeq *$ shows that $\Map(K,qX) \iso \Map(K,X)$. So $q$ is right adjoint to the fully faithful inclusion $\ker(L) \hookrightarrow \D(A)$, and hence $q$ exhibits $\ker(L)$ as the localization of $\D(A)$ at the morphisms inverted by $q$ \cite[Proposition~5.2.7.12]{lurie2009HTT}. Since $q$ is exact, these are the morphisms whose cofiber lies in $\ker(q)$, which is $\D_{(A/I)\text{-}\Mod}(A)$ by \Cref{lem:Almost_Essential_Image}. Hence $\ker(L) \simeq \D(A)/\D_{(A/I)\text{-}\Mod}(A)$.

		By \Cref{prop:Kernel_Homological_Epi} and \Cref{lem:Almost_Restriction_FF}, $\ker(L)$ is dualizable and its inclusion into $\D(A)$ is strongly continuous; explicitly, its right adjoint is $X \mapsto I \otimes^{\bL}_A X$ by \Cref{lem:Almost_Essential_Image}. To see that $\ker(L)$ is also the kernel in $\Pr^{\dual}$, let $E$ be dualizable and $F\colon E \to \D(A)$ strongly continuous with $LF \simeq 0$. Then $F$ factors uniquely through a colimit-preserving functor $F'\colon E \to \ker(L)$, and $F'$ is strongly continuous since its right adjoint is the composite of the colimit-preserving functors $\ker(L) \hookrightarrow \D(A) \xrightarrow{F^R} E$.
	\end{proof}

	\begin{remark}
		\label[remark]{rem:Almost_General}
		In the talk, the last part of \Cref{prop:Almost_Derived_Category} was deduced from the general fact that the inclusion $\Pr^{\dual} \hookrightarrow \PrL_{\st}$ preserves pullbacks along localizations (\Cref{prop:Dualizable_Pullbacks}(2)). Efimov \cite[Example~2.26]{Efimov_Large_Categories} shows that $\D((A,I)^a\text{-}\Mod)$ is dualizable even without any flatness assumption on $I$: in general, $\D(A) \to \D_{(A/I)\text{-}\Mod}(A)$ is still a strongly continuous localization, identified with $B \otimes_A -$ for a connective DG ring $B$ with $H_0(B) = A/I$, and $B = A/I$ if $I$ is flat as a right or left $A$-module.
	\end{remark}

	The following is a special case of Keller's example (\Cref{ex:Keller}), and shows that almost module categories are typically very far from being compactly generated.

	\begin{proposition}
		\label[proposition]{prop:Almost_No_Compacts}
		If $I \subseteq \mathrm{Jac}(A)$ is contained in the Jacobson radical of $A$, then $\D((A,I)^a\text{-}\Mod)$ has no nonzero compact objects.
	\end{proposition}
	\prooffor{prop:Almost_No_Compacts}
	\begin{proof}
		By \Cref{prop:Almost_Derived_Category}, the inclusion $\ker(L) \hookrightarrow \D(A)$ is strongly continuous and hence preserves compact objects. So it suffices to show: if $P \in \Perf(A)$ satisfies $A/I \otimes^{\bL}_A P \simeq 0$, then $P \simeq 0$.

		Suppose $P \neq 0$ and let $n$ be minimal with $H_n(P) \neq 0$. We may represent $P$ by a bounded complex of finitely generated projective modules concentrated in homological degrees $\geq n$: starting from any bounded complex $P_\bullet$ of finitely generated projectives, if $a < n$ is the lowest degree with $P_a \neq 0$, then $P_{a+1} \to P_a$ is surjective as $H_a(P) = 0$, hence split, and we may replace $P_{a+1}$ by its kernel (again finitely generated projective) and remove $P_a$. In this representation, $H_n(P) = \coker(P_{n+1} \to P_n)$ is a finitely generated $A$-module, and by right exactness of the tensor product
		\[
			0 = H_n(A/I \otimes^{\bL}_A P) \cong A/I \otimes_A H_n(P) \cong H_n(P)/IH_n(P).
		\]
		So $H_n(P) = IH_n(P) \subseteq \mathrm{Jac}(A)H_n(P)$, and Nakayama's lemma implies $H_n(P) = 0$, a contradiction.
	\end{proof}

	\begin{example}
		\label[example]{ex:Almost_Modules_Over_Cp}
		In the situation of \Cref{ex:Almost_Modules}(2), $A = \Oo_{\C_p}$ is local with maximal ideal $\mm = \mathrm{Jac}(A)$. Hence $\D((\Oo_{\C_p},\mm)^a\text{-}\Mod)$ is a dualizable category without nonzero compact objects.
	\end{example}

	\subsection*{Hilbert-space models for operator ideals}

	% Source: Clausen, Lectures 2--3 (18.--19.07.2023); the proof below was added when preparing these notes, with AI assistance.
	\begin{proposition}
		\label[proposition]{prop:Operator_Ideal_Models}
		Let $A = B(\Hh)$ for a separable infinite-dimensional complex Hilbert space, and let $J$ be either $\Kk(\Hh)$ or $S_{0+} = \bigcap_{p>0} S_p(\Hh)$. Let $E$ be the additive category of separable complex Hilbert spaces and bounded linear maps, with the split exact structure, and put $D = \Ind(\D^b(E))$. Then $A \to A/J$ is a homological epimorphism. Under Morita theory, its kernel identifies with the full stable subcategory $C_J \subseteq D$ generated under colimits by sequences of Hilbert spaces whose transition maps belong to $J$. Here membership in $J$ means compactness, respectively membership in every Schatten class, also for operators between different Hilbert spaces.

		The completed Hilbert tensor product restricts to a symmetric monoidal structure on $C_J$ whose unit and multiplication functors are strongly continuous. In particular, $C_J \in \CAlg(\Pr^{\dual})$.
	\end{proposition}
	\prooffor{prop:Operator_Ideal_Models}
	\begin{proof}
		We use right $A$-modules in this proof. Every separable Hilbert space is a retract of $\Hh$, and every bounded idempotent has closed image. Hence $V \mapsto \Hom_E(\Hh,V)$ identifies $E$ with the finitely generated projective right $A$-modules, and $D$ with their unbounded derived category. To recover our left-module convention, choose a conjugation $c$ on $\Hh$: the complex-linear transpose $a \mapsto ca^*c$ identifies $A\catop$ with $A$ and preserves both ideals. The intrinsic category $C_J$ itself does not require this choice.

		We first record a factorization fact: given countably many $a_n \in J$, there is a positive injective $d \in J$ with $a_n \in dA$ for every $n$. Include a fixed positive injective operator in $J$ among the $a_n$, and choose $\lambda_n > 0$ sufficiently small that
		\[
			s = \sum_n \lambda_n a_na_n^*
		\]
		converges in operator norm and, when $J = S_{0+}$, in every Schatten quasi-norm $\|{-}\|_{1/m}$, $m \geq 1$. Such a choice is possible by imposing the first $n$ quasi-norm bounds on the $n$th summand, with summable bounds in the corresponding quasi-norm metrics. Then $s$ is positive and injective and belongs to $J$, as does $d=s^{1/2}$. Since $a_na_n^* \leq \lambda_n^{-1}s$, the map $dx \mapsto a_n^*x$ is bounded on the dense image of $d$. Extending it to $\Hh$ and taking adjoints gives $a_n = db_n$ with $b_n \in A$.

		Consequently, $J$ is the directed union of the right ideals $dA$ with $d \in J$ positive and injective. Each $dA$ is a free right module of rank one, so $J$ is flat on the right. Also $J^2=J$: polar decomposition writes $a=u|a|^{1/2}|a|^{1/2}$ with both factors in $J$. Both ideals are closed under positive square roots; for $S_{0+}$ this follows directly from the definition by singular values. The Tor calculation in \Cref{lem:Almost_Restriction_FF} now shows that $A \to A/J$ is a homological epimorphism.

		For a positive injective $d \in J$, put
		\[
			T_d = \bigcup_{n\geq 0} d^{1/2^n}A \subseteq J.
		\]
		Under the identifications $d^{1/2^n}A \cong A$, this is the sequential colimit of free modules with transition maps given by left multiplication by $d^{1/2^{n+1}} \in J$. The family of submodules $T_d$ is directed: apply the factorization fact to all the roots of two given operators to obtain $e$ with both submodules contained in $eA \subseteq T_e$. It covers $J$, so $J$ is a filtered colimit of such telescopes.

		Let $L = - \otimes_A^{\bL} A/J$, and let $C = \ker(L)$. Its inclusion $i\colon C \to D$ has the colimit-preserving right adjoint $r(M)=M\otimes_A^{\bL}J$. This follows by tensoring the fiber sequence $J \to A \to A/J$ with $M$ and using that $L$ is a localization. Since $D$ is generated as a stable category under colimits by $A$, the category $C$ is generated in the same sense by $r(A)=J$. Every $J$-transition telescope lies in $C$, because applying $L$ gives a sequence with zero transition maps, whose colimit is zero. Conversely, the preceding expression for $J$ as a filtered colimit of the $T_d$ puts $J$ in the stable subcategory generated by these telescopes. Thus $C=C_J$.

		The completed Hilbert tensor product extends to a colimit-preserving multiplication $\mu_D\colon D\otimes D\to D$. It preserves compact objects: it sends pairs of Hilbert spaces to Hilbert spaces, and compact objects are obtained from these by finite stable operations and retracts. Thus $\mu_D^R$ preserves colimits. Tensor products of two $J$-transition telescopes again have $J$-transitions, using the cofinal diagonal in $\N\times\N$. For $S_{0+}$, the relevant assertion follows from
		\[
			\sum_{i,j}(s_i(a)s_j(b))^p
			=\left(\sum_i s_i(a)^p\right)\left(\sum_j s_j(b)^p\right)<\infty
		\]
		for every $p>0$; for compact operators it follows by finite-rank approximation. Hence the tensor product restricts to $C_J$. The unit $\C$ belongs to $C_J$, since its identity has finite rank, and is compact there because it is compact in $D$ and the inclusion preserves colimits.

		Finally, $i\otimes i$ has right adjoint $r\otimes r$: the adjunction $i\dashv r$ is an adjunction between colimit-preserving functors, so it can be tensored. The restricted multiplication has right adjoint
		\[
			\mu_{C_J}^R = (r\otimes r)\mu_D^R i,
		\]
		as follows by applying these adjunctions to mapping spectra. All three functors preserve colimits. The unit functor is strongly continuous by compactness of $\C$, and $C_J$ is dualizable by \Cref{prop:Kernel_Homological_Epi}. This proves the last assertion.
	\end{proof}

	\section{Continuous posets}
	\label[section]{sec:Continuous_Posets}

	% Source: Talk 3 (Max Glaser, 27.10.2026); proofs completed with generative AI (Claude)
	We now study a class of non-stable compactly assembled categories in detail, which was already briefly discussed in \Cref{ex:Continuous_Posets}. They will reappear in \Cref{chap:Sheaves}, where we study sheaves on continuous posets. See \cite{Gierz_Continuous_Lattices} for a comprehensive treatment.

	Throughout, let $P$ be a partially ordered set, regarded as a category.

	\begin{lemma}
		\label[lemma]{lem:Ind_Of_Poset}
		The category $\Ind(P)$ is (equivalent to) a partially ordered set. Explicitly, every object of $\Ind(P)$ is of the form $\quotcolim_{s \in S} s$ for a directed subset $S \subseteq P$, and for directed subsets $S,T \subseteq P$ we have
		\begin{align*}
			\Map_{\Ind(P)}\Big(\quotcolim_{s \in S} s, \quotcolim_{t \in T} t\Big) &= \lim_{s \in S\catop}\colim_{t \in T}\Map_P(s,t) \\
			&\simeq \begin{cases} * & \text{if for all } s \in S \text{ there is } t \in T \text{ with } s \leq t, \\ \emptyset & \text{otherwise.}\end{cases}
		\end{align*}
		Hence $\Ind(P)$ is equivalent to the poset of \emph{ideals} of $P$, i.e.\ non-empty, directed, downward closed subsets $J \subseteq P$, ordered by inclusion, via $S \mapsto {\downarrow}S := \{p \in P \mid p \leq s \text{ for some } s \in S\}$.
	\end{lemma}
	\prooffor{lem:Ind_Of_Poset}
	\begin{proof}
		By \Cref{thm:Filtered_Posets_Cofinal}, every object of $\Ind(P)$ is a colimit of representables indexed by a filtered poset $I$, i.e.\ of the form $\quotcolim_{i \in I} p_i$ for a functor $p\colon I \to P$. The image $S = p(I)$ is a directed subset, and the map $I \to S$ is cofinal by \Cref{lem:Cofinality_Posets} (for $s \in S$, the set $\{i \mid s \leq p_i\}$ is non-empty and upward closed in the filtered poset $I$, hence filtered). This proves the first claim. The formula for mapping spaces follows from \Cref{constr:Ind_Large}: each $\Map_P(s,t)$ is either empty or contractible, filtered colimits of such spaces are again empty or contractible (and contractible if and only if one of the terms is), and a limit of such spaces is contractible if and only if all terms are. In particular, all mapping spaces in $\Ind(P)$ are empty or contractible, so $\Ind(P)$ is a poset. Finally, the condition in the formula only depends on ${\downarrow}S$ and ${\downarrow}T$: it says precisely that ${\downarrow}S \subseteq {\downarrow}T$.
	\end{proof}

	Suppose now that $P$ admits directed suprema. Then the colimit functor $\colim\colon \Ind(P) \to P$ sends a directed subset $S$ to $\sup S$, and the Yoneda embedding sends $b$ to ${\downarrow}b$.

	\begin{lemma}
		\label[lemma]{lem:Posets_Accessible}
		Let $P$ be a small partially ordered set admitting directed suprema. Then $P$ is accessible.
	\end{lemma}
	\prooffor{lem:Posets_Accessible}
	\begin{proof}
		Let $\kappa$ be a regular cardinal with $\kappa > |P|$. Let $D \subseteq P$ be a $\kappa$-directed subset. Since $|D| < \kappa$, the set $D$ itself has an upper bound in $D$ (\Cref{lem:Kappa_Filtered_Posets}), i.e.\ $D$ has a largest element. It follows that every $\kappa$-filtered diagram in $P$ has a cofinal constant subdiagram, so that $\kappa$-filtered colimits in $P$ are ``trivial'': every object of $P$ is $\kappa$-compact, and the Yoneda embedding $P \to \Ind_{\kappa}(P)$ is an equivalence, as every $\kappa$-filtered colimit of representable presheaves on $P$ is representable. Hence $P \simeq \Ind_{\kappa}(P^{\kappa})$ is $\kappa$-accessible.
	\end{proof}

	\begin{definition}
		\label[definition]{def:Continuous_Poset}
		A partially ordered set $P$ is called \emph{continuous} if it admits directed suprema and the functor $\colim = \sup\colon \Ind(P) \to P$ admits a left adjoint $\hat{y}\colon P \to \Ind(P)$.
	\end{definition}

	By \Cref{lem:Posets_Accessible}, a small poset is continuous if and only if it is compactly assembled. Unwinding the definitions using \Cref{lem:Ind_Of_Poset}, the existence of the left adjoint $\hat{y}$ means: for every $b \in P$ there is a directed subset $\hat{y}(b) \subseteq P$ such that for every directed subset $T \subseteq P$,
	\begin{equation}
		\label{eq:Yhat_Poset}
		b \leq \sup T \iff \text{for all } s \in \hat{y}(b) \text{ there exists } t \in T \text{ with } s \leq t.
	\end{equation}

	\begin{definition}
		\label[definition]{def:Way_Below}
		Let $P$ be a poset admitting directed suprema, and let $a,b \in P$. We say that $a$ is \emph{way below} $b$, written $a \ll b$, if for every directed subset $T \subseteq P$ with $b \leq \sup T$ there exists $t \in T$ with $a \leq t$. We write
		\[
			P_{\ll b} := \{a \in P \mid a \ll b\}.
		\]
	\end{definition}

	\begin{remark}
		\label[remark]{rem:Way_Below_Basics}
		\leavevmode
		\begin{enumerate}
			\item If $a \ll b$, then $a \leq b$: apply the definition to $T = \{b\}$.
			\item If $a' \leq a \ll b \leq b'$, then $a' \ll b'$.
			\item An element $b$ is compact in $P$ (i.e.\ $\Map_P(b,-)$ preserves filtered colimits) if and only if $b \ll b$.
			\item If $P$ is continuous, then $\hat{y}(b) \subseteq P_{\ll b}$ for every $b$. Indeed, let $s \in \hat{y}(b)$ and let $T$ be directed with $b \leq \sup T$. By \eqref{eq:Yhat_Poset}, there exists $t \in T$ with $s \leq t$, so $s \ll b$.
		\end{enumerate}
	\end{remark}

	\begin{proposition}
		\label[proposition]{prop:Continuous_Poset_Characterization}
		Let $P$ be a poset admitting directed suprema. Then $P$ is continuous if and only if for every $b \in P$ the subset $P_{\ll b}$ is directed and $\sup P_{\ll b} = b$. In this case, $\hat{y}(b) = P_{\ll b}$, viewed as an object of $\Ind(P)$.
	\end{proposition}
	\prooffor{prop:Continuous_Poset_Characterization}
	\begin{proof}
		Suppose that each $P_{\ll b}$ is directed with supremum $b$. We claim that $b \mapsto P_{\ll b}$ satisfies \eqref{eq:Yhat_Poset}, and hence defines a left adjoint to $\sup$ (a monotone map, by \Cref{rem:Way_Below_Basics}(2)). If $b \leq \sup T$, then by definition every $a \ll b$ satisfies $a \leq t$ for some $t \in T$. Conversely, if every $a \in P_{\ll b}$ is bounded by some element of $T$, then $b = \sup P_{\ll b} \leq \sup T$.

		Conversely, suppose that $P$ is continuous. We claim that
		\[
			P_{\ll b} = \{a \in P \mid a \leq s \text{ for some } s \in \hat{y}(b)\} = {\downarrow}\hat{y}(b).
		\]
		For ``$\supseteq$'', note that $\hat{y}(b) \subseteq P_{\ll b}$ by \Cref{rem:Way_Below_Basics}(4), and $P_{\ll b}$ is downward closed by \Cref{rem:Way_Below_Basics}(2). For ``$\subseteq$'', note that $\hat{y}$ is fully faithful (\Cref{lem:Adjoint_Triples}), so the counit $\sup\hat{y}(b) \leq b$ is an equality, i.e.\ $b = \sup \hat{y}(b)$. Applying the definition of $a \ll b$ to the directed set $T = \hat{y}(b)$ gives $s \in \hat{y}(b)$ with $a \leq s$. Since $\hat{y}(b)$ is directed, so is ${\downarrow}\hat{y}(b) = P_{\ll b}$, and $\sup P_{\ll b} = \sup \hat{y}(b) = b$. Finally, $P_{\ll b}$ and $\hat{y}(b)$ define the same object of $\Ind(P)$ by \Cref{lem:Ind_Of_Poset}.
	\end{proof}

	\begin{example}
		\label[example]{ex:Continuous_Poset_Reals}
		The poset $P = \R \cup \{\infty\}$ (with its usual order) admits directed suprema: the supremum of a non-empty subset is its usual supremum if the subset is bounded, and $\infty$ otherwise. We claim that $a \ll b$ if and only if $a < b$. If $a < b$ and $T$ is directed with $b \leq \sup T$, then $\sup T > a$, so some $t \in T$ satisfies $t > a$. Conversely, suppose $a \ll b$; then $a \leq b$ by \Cref{rem:Way_Below_Basics}(1). If $a = b$, consider the directed set $T = \{b - 1/n \mid n \geq 1\}$ (for $b \in \R$) resp.\ $T = \N$ (for $b = \infty$): it has supremum $b$, but no element $t \in T$ satisfies $b \leq t$. Hence $a < b$. It follows that $P_{\ll b} = \{a \mid a < b\}$ is directed with supremum $b$, so $P$ is continuous by \Cref{prop:Continuous_Poset_Characterization}, with $\hat{y}(b) = \quotcolim_{a < b} a$. Note that $P$ has no compact elements.
	\end{example}

	\begin{example}
		\label[example]{ex:Continuous_Poset_Opens}
		Let $X$ be a locally compact Hausdorff space and let $\Open(X)$ be the poset of open subsets of $X$, ordered by inclusion. Directed suprema in $\Open(X)$ are unions. We claim that
		\[
			U \ll V \iff \overline{U} \text{ is compact and } \overline{U} \subseteq V,
		\]
		in which case we write $U \Subset V$. Suppose first that $\overline{U}$ is compact and $\overline{U} \subseteq V$, and let $(W_i)_{i \in I}$ be a directed family of open subsets with $V \subseteq \bigcup_i W_i$. Then the $W_i$ cover the compact set $\overline{U}$, so finitely many of them do, and since the family is directed, a single $W_i$ contains $\overline{U} \supseteq U$. Conversely, let $U \ll V$. Since $X$ is locally compact Hausdorff, every point $x \in V$ has an open neighbourhood $W$ with $\overline{W}$ compact and $\overline{W} \subseteq V$. Hence $V$ is the union of the family $\{W \in \Open(X) \mid W \Subset V\}$, which is directed since $W_1 \cup W_2 \Subset V$ if $W_1,W_2 \Subset V$. By assumption, $U \subseteq W$ for some $W \Subset V$, and then $\overline{U} \subseteq \overline{W}$ is a closed subset of a compact set, hence compact, and $\overline{U} \subseteq \overline{W} \subseteq V$.

		The argument shows moreover that $\Open(X)_{\ll V}$ is directed with union $V$, so $\Open(X)$ is a continuous poset by \Cref{prop:Continuous_Poset_Characterization}, with $\hat{y}(V) = \quotcolim_{U \Subset V} U$. Its compact elements are the open subsets $U$ with $U \Subset U$, i.e.\ the compact open subsets of $X$.
	\end{example}

	\begin{remark}
		\label[remark]{rem:Continuous_Posets_Sheaves}
		If $X$ is a locally compact Hausdorff space, then the morphism $h_U \to h_V$ of representable sheaves on $X$ is a compact morphism in $\Shv(X;\Sp)$ whenever $U \Subset V$ (Mathew, Lecture 3). Since every open $V$ is the union of a sequence $U_0 \Subset U_1 \Subset \cdots$ if $X$ is second countable, this suggests (and, suitably generalized, proves) that $\Shv(X;\Sp)$ is generated by sequential colimits along compact morphisms and hence dualizable, cf.\ \Cref{rem:Omega_1_Compact_Compact_Morphisms}. We will return to this in \Cref{chap:Sheaves}.
	\end{remark}

	\section{Extensions and semiorthogonal gluing}
	\label[section]{sec:Extensions_And_Gluing}

	We first show that semiorthogonal decompositions and gluings behave well with respect to dualizability. We then explain why a general extension behaves differently. The Tate-object classification at the end of this section is a more specialized application and may be omitted on a first reading.

	% Source: Talk 6 (Jacopo Mazzei, 17.11.2026); Efimov, Propositions 2.76 and 2.77
	\begin{proposition}[{\cite[Proposition~2.76]{Efimov_Large_Categories}}]
		\label[proposition]{prop:SOD_Dualizable}
		Let $E = \lra{C,D}$ be a semiorthogonal decomposition of presentable stable categories in the sense of \Cref{def:Semiorthogonal_Decomposition}. Then $E$ is dualizable if and only if $C$ and $D$ are dualizable. In this case $E^{\omega} = \lra{C^{\omega},D^{\omega}}$.
	\end{proposition}
	\prooffor{prop:SOD_Dualizable}
	\begin{proof}
		The projections $i_C^L$ and $i_D^R$ preserve colimits, so $C$ and $D$ are retracts of $E$ in $\PrL_{\st}$. Thus dualizability of $E$ implies that of $C$ and $D$.

		Conversely, if $C$ and $D$ are dualizable, then $E$ is dualizable by \Cref{prop:Generation_Dualizable}, since the strongly continuous inclusions $i_C$ and $i_D$ jointly generate $E$.

		The inclusions preserve compact objects. The projection $i_C^L$ does too, since its right adjoint $i_C$ preserves colimits. For compact $e \in E$, the cofiber sequence $i_Di_D^Re \to e \to i_Ci_C^Le$ therefore shows that $i_Di_D^Re$ is compact. Since $i_D$ is fully faithful and preserves filtered colimits, $i_D^Re$ is compact in $D$. This restricts the semiorthogonal decomposition to compact objects.
	\end{proof}

	\begin{construction}[Gluing along a functor]
		\label[construction]{constr:Semiorthogonal_Gluing}
		Let $F\colon D \to C$ be an accessible exact functor between presentable stable categories. Let $E$ be the category of triples $(x,y,u)$ with $x \in C$, $y \in D$, and $u\colon x \to F(y)$. The same comma category occurs in the construction of a recollement in \href{https://animarium.org/tag/026K/}{\emph{Lectures on Topos Theory}, Construction~5.48}. This is a presentable stable category \cite[Proposition~1.2]{Efimov_Large_Categories}; colimits are computed on $x$ and $y$, using the natural maps $\colim F(y_j) \to F(\colim y_j)$.

		The fully faithful colimit-preserving functors
		\[
			i_C(c) = (c,0,0), \qquad i_D(d) = (0,d,0)
		\]
		have right adjoints $i_C^R(x,y,u) = \fib(u)$ and $i_D^R(x,y,u) = y$. Moreover, $i_C^L(x,y,u) = x$ is left adjoint to $i_C$, and every object fits into a cofiber sequence
		\[
			i_D(y) \to (x,y,u) \to i_C(x).
		\]
		We have $\hom_E(i_D(d),i_C(c)) \simeq 0$. The functor $i_D$ is always strongly continuous, since its right adjoint $(x,y,u) \mapsto y$ preserves colimits.

		In \cite[Definition~1.1]{Efimov_Large_Categories}, the inclusion of $C$ is instead defined as $c \mapsto (c[-1],0,0)$. Its right adjoint is then $(x,y,u) \mapsto \cofib(u)$, and its composite with $i_D$ is $F$ itself, while in our convention $i_C^Ri_D \simeq \Omega F$. This shift plays no role in what follows.
	\end{construction}

	% Source: Talk 6 (Jacopo Mazzei, 17.11.2026); Efimov, Proposition 2.77
	\begin{proposition}[{\cite[Proposition~2.77]{Efimov_Large_Categories}}]
		\label[proposition]{prop:Gluing_Dualizable}
		Let $F\colon D \to C$ be an accessible exact functor between dualizable categories, and let $E$ be the gluing of \Cref{constr:Semiorthogonal_Gluing}. Then the following are equivalent:
		\begin{enumerate}
			\item $F$ preserves colimits;
			\item $E$ is dualizable, and $i_C$ and $i_D$ are strongly continuous.
		\end{enumerate}
		In this case, $i_C$ and $i_D$ exhibit a semiorthogonal decomposition $E = \lra{C,D}$.
	\end{proposition}
	\prooffor{prop:Gluing_Dualizable}
	\begin{proof}
		(2) $\implies$ (1): If $i_C$ is strongly continuous, then $i_C^R$ preserves colimits, and hence so does $i_C^Ri_D \simeq \Omega F$. So $F$ preserves colimits. This does not use the dualizability of $E$.

		(1) $\implies$ (2): If $F$ preserves colimits, then $i_C^R(x,y,u) = \fib(u\colon x \to F(y))$ preserves colimits, so $i_C$ is strongly continuous. The functor $i_D$ is always strongly continuous. The functors $i_C$ and $i_D$ are fully faithful, we have $\hom_E(i_D(d),i_C(c)) \simeq 0$, and every object of $E$ is an extension of an object of $i_C(C)$ by one of $i_D(D)$. So we obtain a semiorthogonal decomposition $E = \lra{C,D}$ in the sense of \Cref{def:Semiorthogonal_Decomposition}, and $E$ is dualizable by \Cref{prop:SOD_Dualizable}.
	\end{proof}

	The strong continuity condition in \Cref{prop:Gluing_Dualizable} is essential: a general strongly continuous localization sequence with dualizable outer terms need not have dualizable middle term (\Cref{ex:Extension_Not_Dualizable}).

	% Source: Talk 6 (Jacopo Mazzei, 17.11.2026); Efimov, Proposition 2.78; proof completed with generative AI (Claude)
	Semiorthogonal decompositions are inherited by categories of strongly continuous functors. For presentable stable categories $T$ and $E$, we write $\Fun^{\mathrm{LL}}(T,E) \subseteq \Fun^{\colim}(T,E)$ for the full subcategory of strongly continuous functors. It is a stable subcategory: shifts of strongly continuous functors are strongly continuous, and the right adjoint of the cofiber of a map $G \to G'$ of colimit-preserving functors is the fiber of the induced map $G'^R \to G^R$ of right adjoints, which preserves colimits if $G^R$ and $G'^R$ do. In the following, we use the notion of a semiorthogonal decomposition of a (not necessarily presentable) stable category from \Cref{rem:SOD_Small}.

	\begin{proposition}[{\cite[Proposition~2.78]{Efimov_Large_Categories}}]
		\label[proposition]{prop:SOD_Functor_Categories}
		Let $E = \lra{i_1(C),i_2(D)}$ be a semiorthogonal decomposition of a presentable stable category, let $\pi_1 := i_1^L\colon E \to C$ and $\pi_2 := i_2^R\colon E \to D$ be the projections, and let $T$ be a presentable stable category.
		\begin{enumerate}
			\item A colimit-preserving functor $G\colon T \to E$ is strongly continuous if and only if $\pi_1G$ and $\pi_2G$ are strongly continuous. Postcomposition with $i_1$ and $i_2$ gives a semiorthogonal decomposition
			\[
				\Fun^{\mathrm{LL}}(T,E) = \lra{\Fun^{\mathrm{LL}}(T,C),\Fun^{\mathrm{LL}}(T,D)}.
			\]
			\item A colimit-preserving functor $G\colon E \to T$ is strongly continuous if and only if $Gi_1$ and $Gi_2$ are strongly continuous. Precomposition with $\pi_1$ and $\pi_2$ gives a semiorthogonal decomposition
			\[
				\Fun^{\mathrm{LL}}(E,T) = \lra{\Fun^{\mathrm{LL}}(C,T),\Fun^{\mathrm{LL}}(D,T)}.
			\]
		\end{enumerate}
	\end{proposition}
	\prooffor{prop:SOD_Functor_Categories}
	\begin{proof}
		Recall from \Cref{prop:SOD_Characterization} that $i_1$, $i_2$, $\pi_1$ and $\pi_2$ are strongly continuous, that $\pi_1i_1 \simeq \id$, $\pi_2i_2 \simeq \id$ and $\pi_1i_2 \simeq 0$, and that there is a cofiber sequence $i_2\pi_2 \to \id_E \to i_1\pi_1$. In particular, composition with any of these four functors preserves strong continuity.

		(1) If $G$ is strongly continuous, then so are $\pi_1G$ and $\pi_2G$. Conversely, if $\pi_1G$ and $\pi_2G$ are strongly continuous, then so are $i_1\pi_1G$ and $i_2\pi_2G$, and hence so is $G$, by the cofiber sequence $i_2\pi_2G \to G \to i_1\pi_1G$ and the stability of $\Fun^{\mathrm{LL}}(T,E)$. Let $j_1 := i_1 \circ -$ and $j_2 := i_2 \circ -$. The pointwise adjunctions $\pi_1 \dashv i_1$ and $i_2 \dashv \pi_2$ show that $\pi_1 \circ -$ is left adjoint to $j_1$ and that $\pi_2 \circ -$ is right adjoint to $j_2$. Hence $j_1$ and $j_2$ are fully faithful, since $\pi_1i_1 \simeq \id$ and $\pi_2i_2 \simeq \id$, and
		\[
			\Map(j_2(H'),j_1(H)) \simeq \Map(\pi_1i_2H',H) \simeq *
		\]
		for $H \in \Fun^{\mathrm{LL}}(T,C)$ and $H' \in \Fun^{\mathrm{LL}}(T,D)$. Finally, every $G$ is an extension $i_2\pi_2G \to G \to i_1\pi_1G$ of an object in the image of $j_1$ by one in the image of $j_2$.

		(2) The criterion for strong continuity follows from \Cref{lem:Automatic_Strong_Continuity}, applied to the functors $i_1$ and $i_2$, whose essential images generate $E$; the converse holds since $i_1$ and $i_2$ are strongly continuous. Let $j_1' := - \circ \pi_1$ and $j_2' := - \circ \pi_2$. For an adjunction $L \dashv R$, precomposition with $R$ is left adjoint to precomposition with $L$. Hence $j_1'$ has the left adjoint $- \circ i_1$, and $j_2'$ has the right adjoint $- \circ i_2$. Again $j_1'$ and $j_2'$ are fully faithful, as $\pi_1i_1 \simeq \id$ and $\pi_2i_2 \simeq \id$. Moreover,
		\[
			\Map(j_2'(H'),j_1'(H)) \simeq \Map(H',H\pi_1i_2) \simeq *
		\]
		for $H \in \Fun^{\mathrm{LL}}(C,T)$ and $H' \in \Fun^{\mathrm{LL}}(D,T)$. Finally, precomposing $G$ with the cofiber sequence $i_2\pi_2 \to \id_E \to i_1\pi_1$ gives a cofiber sequence $j_2'(Gi_2) \to G \to j_1'(Gi_1)$.
	\end{proof}

	\begin{warning}
		\label[warning]{warn:Extensions_Not_Dualizable}
		The converse of \Cref{prop:Localization_Sequence_Dualizable} fails: if $K$ and $E$ are dualizable, then $D$ need not be dualizable, even if $K$ and $E$ are compactly generated. An explicit counterexample is given in \Cref{ex:Extension_Not_Dualizable} below. However, dualizable categories are closed under pullbacks along strongly continuous localizations, see \Cref{prop:Dualizable_Pullbacks}(2).
	\end{warning}

	% Source: Efimov, Lecture 4 (MPIM Bonn, 10.09.2024)
	Let us take a closer look at extensions of compactly generated categories. Consider a strongly continuous localization sequence
	\[
		\Ind(A) \xrightarrow{F} D \xrightarrow{G} \Ind(B)
	\]
	in $\PrL_{\st}$, with $A,B \in \Cat_{\perf}$. Since $\hom_D(Fa,G^Rb) \simeq \hom_{\Ind(B)}(GFa,b) \simeq 0$, the category $D$ admits a semiorthogonal decomposition
	\[
		D = \lra{G^R(\Ind(B)),F(\Ind(A))}
	\]
	in the ordinary sense of semiorthogonality and generation. It need not satisfy \Cref{def:Semiorthogonal_Decomposition}: although $G^R$ preserves colimits, its right adjoint need not do so. Thus \Cref{prop:SOD_Dualizable} does not apply to an arbitrary extension. Since $G$ is strongly continuous, $G^R$ admits a right adjoint $G^{RR}$, and the extension $D$ is encoded by the accessible functor
	\[
		G^{RR} \circ F\colon \Ind(A) \to \Ind(B).
	\]
	An accessible exact functor $\Phi\colon \Ind(A) \to \Ind(B)$ is in turn determined by the functor $\Psi\colon B \to \Pro(\Ind(A))$ sending $b \in B$ to the pro-object corepresenting the accessible functor $x \mapsto \hom_{\Ind(B)}(b,\Phi(x))$.

	Following Hennion \cite{Hennion_Tate}, we denote by $\mathrm{Tate}^{\mathrm{el}}(A) \subseteq \Pro(\Ind(A))$ the stable subcategory generated by $\Ind(A)$ and $\Pro(A)$ (the \emph{elementary Tate objects}), and by $\mathrm{Tate}(A) \subseteq \Pro(\Ind(A))$ its idempotent completion. By \cite[Theorem~3]{Hennion_Tate}, an object of $\Pro(\Ind(A))$ lies in $\mathrm{Tate}^{\mathrm{el}}(A)$ if and only if it is an extension $u \to t \to v$ with $u \in \Pro(A)$ and $v \in \Ind(A)$.

	\begin{proposition}[{Efimov \cite[Proposition~6.3]{Efimov_Large_Categories}}]
		\label[proposition]{prop:Extension_Dualizable_Tate}
		In the above situation, let $\Psi\colon B \to \Pro(\Ind(A))$ be the functor corresponding to $\Phi = G^{RR} \circ F$. Then the following are equivalent:
		\begin{enumerate}
			\item $D$ is dualizable;
			\item $D$ is compactly generated;
			\item The essential image of $\Psi$ is contained in $\mathrm{Tate}(A)$.
		\end{enumerate}
	\end{proposition}
	\prooffor{prop:Extension_Dualizable_Tate}
	\begin{proof}
		We first show that $D$ is compactly generated if and only if for every $b \in B$ there is a compact object $d \in D^{\omega}$ with $G(d) \simeq b \oplus \Sigma b$. If $D$ is compactly generated, this follows from \Cref{thm:Thomason_Neeman}(2) and \Cref{rem:Lifting_Objects_Quotient}. Conversely, suppose such lifts $d_b$ exist. Since $F$ is strongly continuous, the objects $F(a)$, $a \in A$, are compact, and we claim that together with the $d_b$ they generate $D$. Indeed, if $x \in D$ is right orthogonal to all of them, then $F^R(x) = 0$, so $x \simeq G^RG(x)$, and $0 \simeq \hom_D(d_b, G^RG(x)) \simeq \hom_{\Ind(B)}(b \oplus \Sigma b, G(x))$ for all $b$, so $G(x) = 0$ and hence $x = 0$.

		(1) $\Rightarrow$ (2): If $D$ is dualizable, then the required lifts exist by \cite[Proposition~2.10]{Efimov_Large_Categories} (which reduces to the compactly generated case by embedding $D$ strongly continuously into a compactly generated category). (2) $\Rightarrow$ (1) is clear.

		(2) $\Leftrightarrow$ (3): Let $d \in D$, and put $z := G(d)$ and $w := F^R(d)$, so that we have a fiber sequence $F(w) \to d \to G^R(z)$. Applying $\hom_D(-,F(w'))$ for $w' \in \Ind(A)$ and using $\hom_D(G^R(z),F(w')) \simeq \hom_{\Ind(B)}(z,\Phi(w'))$, we obtain a fiber sequence of functors $\Ind(A) \to \Sp$
		\[
			\hom_{\Ind(B)}(z,\Phi(-)) \to \hom_D(d,F(-)) \to \hom_{\Ind(A)}(w,-).
		\]
		Every filtered colimit in $D$ is an extension of a filtered colimit in $F(\Ind(A))$ by one in $G^R(\Ind(B))$, as $F^R$ and $G$ preserve colimits. Since $\hom_D(d,G^R(-)) \simeq \hom_{\Ind(B)}(z,-)$, it follows that $d$ is compact if and only if $z$ is compact and $\hom_D(d,F(-))$ preserves filtered colimits. Now identify exact accessible functors $\Ind(A) \to \Sp$ with $\Pro(\Ind(A))\catop$; under this identification, the functors preserving filtered colimits correspond to $\Pro(A)$. The first term of the fiber sequence corresponds to $\Psi(z)$ by definition, and the last one to $w \in \Ind(A)$. So $d$ is compact if and only if $z \in B$ and there is a cofiber sequence $u \to \Psi(z) \to \Sigma w$ in $\Pro(\Ind(A))$ with $u \in \Pro(A)$. Conversely, given $z \in B$ and a cofiber sequence $u \to \Psi(z) \to v$ with $u \in \Pro(A)$ and $v \in \Ind(A)$, the map $\Psi(z) \to v$ corresponds to a map $z \to \Phi(v)$, i.e.\ to a map $G^R(z) \to F(v)$, and its fiber $d$ is a compact object with $G(d) \simeq z$ by the same computation. Using Hennion's theorem, we conclude that $z \in B$ lifts to $D^{\omega}$ if and only if $\Psi(z) \in \mathrm{Tate}^{\mathrm{el}}(A)$.

		Combining this with the first paragraph, $D$ is compactly generated if and only if $\Psi(b) \oplus \Sigma\Psi(b) \in \mathrm{Tate}^{\mathrm{el}}(A)$ for all $b \in B$. This implies $\Psi(b) \in \mathrm{Tate}(A)$. Conversely, if $\Psi(b) \in \mathrm{Tate}(A)$, then $\Psi(b) \oplus \Sigma\Psi(b) \in \mathrm{Tate}^{\mathrm{el}}(A)$ by Thomason's classification of dense subcategories, as in \Cref{rem:Lifting_Objects_Quotient}.
	\end{proof}

	\begin{remark}
		\label[remark]{rem:Tate_Classification_Of_Extensions}
		Applied to compactly generated $D$, the proof shows that short exact sequences $A \to D_0 \to B$ in $\Cat_{\perf}$ correspond to exact functors $B \to \mathrm{Tate}(A)$; this classification first appeared in \cite[Theorem~A.5.1 and Corollary~A.5.10]{Hermitian_II}, see also \cite[Theorem~6.4]{Efimov_Large_Categories}.
	\end{remark}

	\begin{example}[Efimov]
		\label[example]{ex:Extension_Not_Dualizable}
		Let $k$ be a field, and consider the category
		\[
			D := \Big\{(V,W,\phi) \;\Big|\; V,W \in \D(k),\ \phi\colon \bigoplus_{\N}V \to \bigoplus_{\N}W\Big\}.
		\]
		The functors $W \mapsto (0,W,0)$ and $(V,W,\phi) \mapsto V$ exhibit $D$ as an extension of $\D(k)$ by $\D(k)$, both of which are compactly generated. Nevertheless, $D$ is not dualizable.

		\emph{Proof} (following \cite[Proposition~6.8]{Efimov_Large_Categories})\emph{.} Here $\Phi(W) = \prod_{\N}\bigoplus_{\N}W$, since $\hom(\bigoplus_{\N}V,\bigoplus_{\N}W) \simeq \hom(V,\prod_{\N}\bigoplus_{\N}W)$. By \Cref{prop:Extension_Dualizable_Tate}, it suffices to show that $\Psi(k) \notin \mathrm{Tate}(\Perf(k))$, where $\Psi(k)$ is the pro-object corepresenting $W \mapsto \hom(k,\Phi(W)) = \prod_{\N}\bigoplus_{\N}W$. Using (AB6) in $\D(k)$, we may rewrite this functor as
		\[
			\prod_{\N}\bigoplus_{\N}W \simeq \colim_{f\colon \N \to \N}\ \prod_{n \in \N}\bigoplus_{i=0}^{f(n)-1}W \simeq \colim_{f\colon \N \to \N}\hom(X_f,W), \qquad X_f := \bigoplus_{n \in \N}k^{f(n)},
		\]
		where $\N^{\N}$ is ordered pointwise, and for $f \leq g$ the transition map is induced by the projection $X_g \to X_f$. Hence $\Psi(k) \simeq \text{``}\!\lim\nolimits_f\!\text{} X_f$, a pro-system in which all transition maps are split epimorphisms.

		Let $q\colon \D(k) \to \Calk(\Perf(k))$ be the projection to $\Calk(\Perf(k)) := \bigcup_{\kappa}\Calk_{\kappa}(\Perf(k))$ (the transition functors are fully faithful by \Cref{rem:Calkin_Mapping_Spectra}); this is an idempotent complete stable category, and $q(x) = 0$ if and only if $x \in \Perf(k)$. Suppose that $\Psi(k) \in \mathrm{Tate}(\Perf(k))$. Then $\Psi(k) \oplus \Sigma\Psi(k) \in \mathrm{Tate}^{\mathrm{el}}(\Perf(k))$, so there is a cofiber sequence $u \to \Psi(k) \oplus \Sigma\Psi(k) \to v$ with $u \in \Pro(\Perf(k))$ and $v \in \D(k)$. Applying $\Pro(q)$ kills $u$, so $\Pro(q)(\Psi(k)) = \text{``}\!\lim\nolimits_f\!\text{} q(X_f)$ is a retract of the constant pro-object $q(v)$, and hence is itself isomorphic to a constant pro-object $c$.

		We now use the following elementary observation: if a pro-system $(Y_f)_f$ indexed by a directed poset, with split epimorphisms as transition maps $t_{gf}\colon Y_g \to Y_f$, is isomorphic to a constant pro-object $c$, then there is $f_0$ such that $t_{gf}$ is an isomorphism for all $f_0 \leq f \leq g$. Indeed, the isomorphism is given by maps $p_f\colon Y_f \to c$ for $f \geq f_0$ (compatible with the $t_{gf}$) and a compatible family $s_f\colon c \to Y_f$, such that $p_fs_f = \id_c$ for $f \geq f_0$, and for every $f$ there is $g \geq f,f_0$ with $s_fp_g = t_{gf}$. If $\sigma$ is a section of $t_{gf}$, then $s_f(p_g\sigma) = \id$. So for $f \geq f_0$, the map $s_f$ is both a split mono and a split epi, hence an isomorphism, and then $t_{gf} = s_fs_g^{-1}$ is an isomorphism for $f_0 \leq f \leq g$.

		Applied to $Y_f = q(X_f)$, this gives a contradiction: for every $f$, the fiber of $X_{f+1} \to X_f$ is $\bigoplus_{\N}k \notin \Perf(k)$, so $q(X_{f+1}) \to q(X_f)$ is not an isomorphism.
	\end{example}

